\Cref{tab:8-threats} consolidates threats to the main interpretations.
The shared-input, timing and provenance controls are defined in Chapter 3;
\cref{sec:7-8} evaluates repeated execution, fixed-point rebuilding and
precision sensitivity within their recorded boundaries.


\par\medskip\noindent\begin{minipage}{\linewidth}
\centering\captionsetup{hypcap=false}
\captionof{table}[Threats to validity]{Threats that can reverse or narrow a conclusion.}
\label{tab:8-threats}
\footnotesize\setlength{\tabcolsep}{4pt}
\begin{tabularx}{\linewidth}{@{}*{3}{>{\raggedright\arraybackslash}X}@{}}
\toprule
Claim at risk & Alternative or reversal & Control and remaining test \\
\midrule
Execution gain & Changed scores, ties or index state explain an apparent
speedup & Strict T4 CPU sorting and table-form diagnostics; finite
coverage, with candidate dumps no longer retained \\
Budget quality & Finite-pool sampling or drift hides damaging queries
& Disjoint GT evaluation, full traces, subset-sum checks and FULL
enumeration; future-query and arbitrary-subset invariance remain untested \\
Competitive margin & Grid density, call tails or API overhead changes
the ordering & Both recall estimators, retained slow calls and int8
baseline; TF32/FP32 sensitivity; T5 exposes small-batch reversals;
cross-device comparisons remain unmeasured \\
Capacity/build claim & Transient allocations exceed a resident sample
& Phase bytes, repeated empty-cache builds and 20\,ms device peaks;
sub-interval allocations and full retrained frontiers remain unmeasured \\
\bottomrule
\end{tabularx}
\end{minipage}\par\medskip

Coverage is one GPU, squared-L2 top-10 and six prepared embedding
datasets (\cref{sec:7-smallbatch}). Positional query removal does not
exclude near-duplicate embeddings, and the exports lack relevance labels.
Updates, concurrent traffic, changing query populations and deployment
latency targets remain untested, as do billion-scale and multi-device
capacity.
